\begin{frame}[t]{Esquema de um botão com pull-up interno}
\begin{center}\begin{tikzpicture}[x=1cm,y=0.9cm,every node/.style={font=\small}]
\draw[dashed,darkgreen] (-3,-2) rectangle (0.9,2.5);
\node at (-1.1,2.1) {Dentro do Uno R3};
\node at (-1,1.4) {5 V};
\draw (-1,1.1)--(-1,0.8);
\draw[fill=darkgreen!10] (-1.2,0.15) rectangle (-0.8,0.8);
\node[left] at (-1.25,0.5) {Pull-up};
\draw (-1,0.15)--(-1,-0.5)--(1.8,-0.5);
\fill (-1,-0.5) circle (1.5pt);
\draw (-1,-0.5)--(-2.1,-0.5);\node[left] at (-2.1,-0.5) {D2};
\draw (1.8,-0.5)--(2.5,0);\draw (2.65,-0.5)--(3.4,-0.5)--(3.4,-1.5);
\fill (1.8,-0.5) circle (1.5pt);\fill (2.65,-0.5) circle (1.5pt);
\node[above] at (2.3,0.1) {Botão};
\draw (3,-1.5)--(3.8,-1.5);\draw (3.1,-1.65)--(3.7,-1.65);\draw (3.2,-1.8)--(3.6,-1.8);
\node[right] at (3.8,-1.5) {GND};
\end{tikzpicture}\end{center}
\small Ligação externa: um contato do botão no D2; o outro no GND. Use os contatos que fecham ao pressionar e confira a orientação em botões de quatro terminais.
\end{frame}
